% CORRECTED VERSION
% Changes:
% - Fixed the inequality in Step 6 by requiring η_max ≤ 1/L and adding Lemma 2.
% - Removed the spurious expectation on the deterministic learning rate in Step 7.
% - Added Lemma 3 to justify the lower bound on the average learning rate for any T,
%   handling the case when T is not a multiple of the restart period M.
% - Clarified the definition of the random index τ in Step 10.
% - Updated Assumption 5 to reflect the correct step‑size bound.

\documentclass{article}
\usepackage{amsmath,amssymb,amsthm}
\usepackage{algorithm}
\usepackage{algorithmic}
\usepackage{enumitem}
\usepackage{hyperref}
\usepackage{bm}
\usepackage{geometry}
\geometry{margin=1in}
\newtheorem{theorem}{Theorem}
\newtheorem{lemma}{Lemma}
\newtheorem{assumption}{Assumption}
\newtheorem{remark}{Remark}
\newcommand{\EE}{\mathbb{E}}
\newcommand{\RR}{\mathbb{R}}
\newcommand{\norm}[1]{\left\|#1\right\|}
\newcommand{\inner}[2]{\left\langle #1,#2\right\rangle}
\begin{document}

\section*{Algorithm Name}
\textbf{Stochastic Gradient Descent with Warm Restarts (SGDR)}  

The method follows standard stochastic gradient descent while the learning rate follows a cosine‑annealing schedule that is periodically reset (warm restart) every $M$ iterations.

\section*{Mathematical Formulation}
Let $f:\RR^{d}\!\to\!\RR$ be the (possibly non‑convex) objective and let $\{z_t\}_{t\ge 0}$ be i.i.d.\ data samples.  
Define the stochastic gradient
\[
g_t \;:=\; \nabla f(w_t;z_t),\qquad 
\EE\!\left[g_t\mid w_t\right]=\nabla f(w_t).
\]

The learning‑rate schedule is
\begin{equation}
\boxed{
\eta_t \;=\; \eta_{\max}\,
\frac{1}{2}\Bigl(1+\cos\!\bigl(\pi\,\frac{t\bmod M}{M}\bigr)\Bigr)
\quad\text{for }t=0,1,\dots
}
\label{eq:eta}
\end{equation}
Thus $\eta_t$ starts from $\eta_{\max}$, decays smoothly to $0$ after $M$ steps, and then is reset to $\eta_{\max}$ (a warm restart).  

Update rule:
\begin{equation}
\boxed{
w_{t+1}=w_t-\eta_t\,g_t,\qquad t=0,1,\dots
}
\label{eq:update}
\end{equation}

Hyper‑parameters: $\eta_{\max}>0$ (max step size) and $M\in\mathbb{N}$ (restart period).

\section*{Pseudocode}
\begin{algorithm}[H]
\caption{SGDR}
\begin{algorithmic}[1]
\STATE \textbf{Input:} $\eta_{\max}>0$, restart period $M$, initial point $w_0$.
\FOR{$t=0,1,\dots,T-1$}
    \STATE Compute stochastic gradient $g_t=\nabla f(w_t;z_t)$.
    \STATE Compute learning rate 
    \[
    \eta_t\gets \eta_{\max}\frac{1}{2}\Bigl(1+\cos\!\bigl(\pi\frac{t\bmod M}{M}\bigr)\Bigr).
    \]
    \STATE Update $w_{t+1}\gets w_t-\eta_t g_t$.
\ENDFOR
\STATE \textbf{Return} $w_T$ (or the average $\bar w_T=\frac1T\sum_{t=0}^{T-1}w_t$).
\end{algorithmic}
\end{algorithm}

\section*{Dry Run Example}
Consider the one‑dimensional quadratic $f(w)=(w-3)^2$, whose gradient is $\nabla f(w)=2(w-3)$.  
Set $w_0=0$, $\eta_{\max}=0.1$, $M=4$ (so we will not see a full restart in the first three steps).

\begin{center}
\begin{tabular}{c|c|c|c|c}
$t$ & $\eta_t$ (from \eqref{eq:eta}) & $g_t=2(w_t-3)$ & $w_{t+1}=w_t-\eta_t g_t$ & $f(w_{t+1})$ \\ \hline
0 & $0.1\frac{1+\cos0}{2}=0.1$ & $-6$ & $0-0.1(-6)=0.6$ & $(0.6-3)^2=5.76$\\
1 & $0.1\frac{1+\cos(\pi/4)}{2}=0.1\cdot0.8536=0.08536$ & $2(0.6-3)=-4.8$ & $0.6-0.08536(-4.8)=1.0097$ & $(1.0097-3)^2\approx3.95$\\
2 & $0.1\frac{1+\cos(\pi/2)}{2}=0.1\cdot0.5=0.05$ & $2(1.0097-3)=-3.9806$ & $1.0097-0.05(-3.9806)=1.2097$ & $(1.2097-3)^2\approx3.20$\\
\end{tabular}
\end{center}
The objective decreases at each iteration, illustrating the effect of the decaying learning rate.

\newpage
\section*{Assumptions}
\begin{assumption}[Smoothness]\label{ass:smooth}
$f$ is $L$‑smooth, i.e.
\[
\forall w,u\in\RR^{d}:\quad 
f(u)\le f(w)+\inner{\nabla f(w)}{u-w}+\frac{L}{2}\norm{u-w}^{2}.
\]
\end{assumption}
\begin{assumption}[Lower Boundedness]\label{ass:lb}
$f$ is bounded below: $f^{\star}:=\inf_{w}f(w)>-\infty$.
\end{assumption}
\begin{assumption}[Unbiased Gradient]\label{ass:unbiased}
For all $t$, $\EE\!\left[g_t\mid w_t\right]=\nabla f(w_t)$.
\end{assumption}
\begin{assumption}[Bounded Variance]\label{ass:variance}
There exists $\sigma^{2}>0$ such that
\[
\EE\!\left[\norm{g_t-\nabla f(w_t)}^{2}\mid w_t\right]\le\sigma^{2},\qquad\forall t.
\]
\end{assumption}
\begin{assumption}[Step‑size Schedule]\label{ass:eta}
The schedule \eqref{eq:eta} satisfies
\[
0\le\eta_t\le\eta_{\max}\quad\forall t,
\qquad
\eta_{\max}\le\frac{1}{L}.
\tag{A5}
\]
Moreover, for any $T\ge M$ we have the average‑learning‑rate lower bound
\[
\frac{1}{T}\sum_{t=0}^{T-1}\eta_t\;\ge\;\frac{\eta_{\max}}{4}.
\tag{A6}
\]
(The constant $1/4$ is convenient for the analysis; a tighter bound $\eta_{\max}/2$ holds when $T$ is an integer multiple of $M$, see Lemma \ref{lem:avg-eta}.) 
\end{assumption}

\section*{Main Convergence Theorem}
\begin{theorem}\label{thm:main}
Assume \ref{ass:smooth}–\ref{ass:eta} hold and run SGDR for $T$ iterations.  
Define the random iterate $w_{\tau}$ where $\tau$ is drawn uniformly from $\{0,\dots,T-1\}$. Then
\[
\EE\!\bigl[\norm{\nabla f(w_{\tau})}^{2}\bigr]
\;\le\;
\frac{2\bigl(f(w_{0})-f^{\star}\bigr)}{\eta_{\max}T}
\;+\;
\frac{L\eta_{\max}\sigma^{2}}{2}.
\tag{1}
\]
Consequently, if $\eta_{\max}=c/\sqrt{T}$ for a constant $c\le 1/L$, the bound becomes
\[
\EE\!\bigl[\norm{\nabla f(w_{\tau})}^{2}\bigr}
\;=\;O\!\Bigl(\frac{1}{\sqrt{T}}\Bigr).
\]
\end{theorem}

\section*{Theoretical Properties}
\begin{itemize}[leftmargin=*]
\item \textbf{Stability.} The cosine annealing never yields a negative step size; the bound \eqref{thm:main} shows that the algorithm remains stable provided $\eta_{\max}\le 1/L$ (so that the descent term dominates the smoothness term).
\item \textbf{Sensitivity to $M$.} The average learning rate over a full cycle equals $\eta_{\max}/2$, independent of $M$. Larger $M$ yields slower decay within a cycle, which empirically improves exploration; the theorem only needs the weaker bound $\eta_{\max}/4$ that holds for any $T\ge M$ (Lemma \ref{lem:avg-eta}).
\item \textbf{Comparison to plain SGD.} For constant step size $\eta$, classic SGD obtains $\EE\!\bigl[\norm{\nabla f(w_{\tau})}^{2}\bigr]\le\frac{2(f(w_{0})-f^{\star})}{\eta T}+L\eta\sigma^{2}$. SGDR replaces $\eta$ by its average $\ge\eta_{\max}/4$, thus enjoys the same $O(1/\sqrt{T})$ rate while periodically resetting the step size, which can help escape shallow local minima.
\end{itemize}

\section*{Discussion}
The theorem shows that warm restarts do not deteriorate the worst‑case convergence guarantee for non‑convex smooth objectives. The bound consists of a deterministic decay term (first term) and a stochastic variance term (second term). By letting $\eta_{\max}=c/\sqrt{T}$ the variance term scales as $O(1/\sqrt{T})$, matching the optimal rate for stochastic first‑order methods under the given assumptions.

\newpage
\section*{Preliminaries (Definitions and Lemmas)}
\begin{lemma}[Descent Lemma]\label{lem:descent}
If $f$ is $L$‑smooth then for any $w$ and any vector $d$,
\[
f(w-d)\le f(w)-\inner{\nabla f(w)}{d}+\frac{L}{2}\norm{d}^{2}.
\]
\end{lemma}
\begin{proof}
Directly from Definition \ref{ass:smooth} with $u=w-d$.
\end{proof}

% ADDED: Lemma establishing the quadratic lower bound used in Step 6.
\begin{lemma}\label{lem:quadratic}
For any $0\le \eta\le 1/L$,
\[
\eta-\frac{L}{2}\eta^{2}\;\ge\;\frac{\eta}{2}.
\]
\end{lemma}
\begin{proof}
Consider the function $h(\eta)=\eta-\frac{L}{2}\eta^{2}-\frac{\eta}{2}
= \frac{\eta}{2}\bigl(1-L\eta\bigr)$.  
Since $0\le\eta\le 1/L$, we have $1-L\eta\ge0$, hence $h(\eta)\ge0$, which is exactly the desired inequality.
\end{proof}

% ADDED: Lemma that gives a uniform lower bound on the average learning rate for any horizon $T$.
\begin{lemma}[Average learning‑rate lower bound]\label{lem:avg-eta}
Let $\eta_t$ be defined by \eqref{eq:eta} with period $M$. For any integer $T\ge M$,
\[
\frac{1}{T}\sum_{t=0}^{T-1}\eta_t\;\ge\;\frac{\eta_{\max}}{4}.
\]
\end{lemma}
\begin{proof}
Write $T = qM+r$ with $0\le r<M$. Over each full period the average equals $\eta_{\max}/2$:
\[
\frac{1}{M}\sum_{t=0}^{M-1}\eta_t
= \frac{\eta_{\max}}{2}.
\]
Hence
\[
\sum_{t=0}^{T-1}\eta_t
= q\sum_{t=0}^{M-1}\eta_t+\sum_{t=0}^{r-1}\eta_t
\ge q\,\frac{\eta_{\max}M}{2},
\]
because every $\eta_t$ is non‑negative. Dividing by $T=qM+r$ yields
\[
\frac{1}{T}\sum_{t=0}^{T-1}\eta_t
\ge \frac{q}{q+r/M}\,\frac{\eta_{\max}}{2}
\ge \frac{1}{2}\,\frac{q}{q+1}\,\eta_{\max}
\ge \frac{\eta_{\max}}{4},
\]
where the last inequality uses $q\ge1$ (i.e. $T\ge M$). ∎
\end{proof}

\section*{Main Proof (Step‑by‑Step Derivation)}
We prove Theorem \ref{thm:main}. Throughout the proof we condition on the filtration 
\[
\mathcal{F}_t:=\sigma(w_0,g_0,\dots,g_{t-1}).
\]

\begin{enumerate}
\item[Step 1.] Apply Lemma \ref{lem:descent} with $d=\eta_t g_t$:
\[
f(w_{t+1})\le f(w_t)-\eta_t\inner{\nabla f(w_t)}{g_t}
+\frac{L\eta_t^{2}}{2}\norm{g_t}^{2}.
\tag{2}
\]

\item[Step 2.] Take expectation conditioned on $\mathcal{F}_t$ and use unbiasedness (Assumption \ref{ass:unbiased}):
\[
\EE\!\bigl[f(w_{t+1})\mid\mathcal{F}_t\bigr]
\le f(w_t)-\eta_t\norm{\nabla f(w_t)}^{2}
+\frac{L\eta_t^{2}}{2}\EE\!\bigl[\norm{g_t}^{2}\mid\mathcal{F}_t\bigr].
\tag{3}
\]

\item[Step 3.] Decompose $\norm{g_t}^{2}$:
\[
\EE\!\bigl[\norm{g_t}^{2}\mid\mathcal{F}_t\bigr]
= \norm{\nabla f(w_t)}^{2}
+ \EE\!\bigl[\norm{g_t-\nabla f(w_t)}^{2}\mid\mathcal{F}_t\bigr]
\le \norm{\nabla f(w_t)}^{2}+\sigma^{2},
\tag{4}
\]
where the inequality follows from Assumption \ref{ass:variance}.

\item[Step 4.] Substitute \eqref{4} into \eqref{3}:
\[
\EE\!\bigl[f(w_{t+1})\mid\mathcal{F}_t\bigr]
\le f(w_t)-\eta_t\norm{\nabla f(w_t)}^{2}
+\frac{L\eta_t^{2}}{2}\bigl(\norm{\nabla f(w_t)}^{2}+\sigma^{2}\bigr).
\tag{5}
\]

\item[Step 5.] Rearrange \eqref{5} to isolate the gradient term:
\[
\bigl(\eta_t-\tfrac{L}{2}\eta_t^{2}\bigr)\norm{\nabla f(w_t)}^{2}
\le f(w_t)-\EE\!\bigl[f(w_{t+1})\mid\mathcal{F}_t\bigr]
+\tfrac{L}{2}\eta_t^{2}\sigma^{2}.
\tag{6}
\]

\item[Step 6.] % CORRECTED: use Lemma \ref{lem:quadratic}
By Lemma \ref{lem:quadratic} (which requires $\eta_t\le 1/L$, guaranteed by Assumption \ref{ass:eta}), we have
\[
\eta_t-\frac{L}{2}\eta_t^{2}\;\ge\;\frac{\eta_t}{2}.
\tag{7}
\]
Thus \eqref{6} yields
\[
\frac{\eta_t}{2}\norm{\nabla f(w_t)}^{2}
\le f(w_t)-\EE\!\bigl[f(w_{t+1})\mid\mathcal{F}_t\bigr]
+\frac{L}{2}\eta_t^{2}\sigma^{2}.
\tag{8}
\]

\item[Step 7.] % CORRECTED: $\eta_t$ is deterministic
Take total expectation and sum over $t=0,\dots,T-1$:
\[
\sum_{t=0}^{T-1}\frac{\eta_t}{2}\,\EE\!\bigl[\norm{\nabla f(w_t)}^{2}\bigr]
\le f(w_0)-\EE\!\bigl[f(w_T)\bigr]
+\frac{L\sigma^{2}}{2}\sum_{t=0}^{T-1}\eta_t^{2}.
\tag{9}
\]

\item[Step 8.] Drop the non‑negative term $-\EE[f(w_T)]\le0$ and use $\eta_t^{2}\le\eta_{\max}^{2}$:
\[
\sum_{t=0}^{T-1}\eta_t\,\EE\!\bigl[\norm{\nabla f(w_t)}^{2}\bigr]
\le
2\bigl(f(w_0)-f^{\star}\bigr)
+L\eta_{\max}^{2}\sigma^{2}T.
\tag{10}
\]

\item[Step 9.] % ADDED justification for the average learning rate
Apply Lemma \ref{lem:avg-eta} (Assumption \ref{ass:eta} (A6)):
\[
\frac{1}{T}\sum_{t=0}^{T-1}\eta_t\;\ge\;\frac{\eta_{\max}}{4}.
\]
Dividing \eqref{10} by $T$ and using the above lower bound gives
\[
\frac{\eta_{\max}}{4}\,
\frac{1}{T}\sum_{t=0}^{T-1}\EE\!\bigl[\norm{\nabla f(w_t)}^{2}\bigr]
\le
\frac{2\bigl(f(w_0)-f^{\star}\bigr)}{T}
+L\eta_{\max}^{2}\sigma^{2}.
\tag{11}
\]

\item[Step 10.] % ADDED clarification of the random index
Let $\tau$ be drawn uniformly from $\{0,\dots,T-1\}$, independent of all randomness in the algorithm. Then
\[
\EE\!\bigl[\norm{\nabla f(w_{\tau})}^{2}\bigr]
= \frac{1}{T}\sum_{t=0}^{T-1}\EE\!\bigl[\norm{\nabla f(w_t)}^{2}\bigr].
\]
Multiplying \eqref{11} by $4/\eta_{\max}$ yields exactly the bound \eqref{1}:
\[
\EE\!\bigl[\norm{\nabla f(w_{\tau})}^{2}\bigr]
\le
\frac{2\bigl(f(w_{0})-f^{\star}\bigr)}{\eta_{\max}T}
+\frac{L\eta_{\max}\sigma^{2}}{2}.
\tag{12}
\]
∎
\end{enumerate}

\section*{Rate Derivation (With Constants)}
Setting $\eta_{\max}=c/\sqrt{T}$ with $c\le 1/L$ gives
\[
\EE\!\bigl[\norm{\nabla f(w_{\tau})}^{2}\bigr]
\le \frac{2\bigl(f(w_{0})-f^{\star}\bigr)}{c}\frac{1}{\sqrt{T}}
+\frac{Lc\sigma^{2}}{2}\frac{1}{\sqrt{T}}
=O\!\Bigl(\frac{1}{\sqrt{T}}\Bigr).
\]

\section*{Special Cases}
\begin{itemize}[leftmargin=*]
\item \textbf{Deterministic case ($\sigma=0$).} The bound reduces to
\[
\EE\!\bigl[\norm{\nabla f(w_{\tau})}^{2}\bigr]\le
\frac{2\bigl(f(w_0)-f^{\star}\bigr)}{\eta_{\max}T},
\]
i.e., a $O(1/T)$ decay of the gradient norm, the standard rate for deterministic gradient descent with diminishing steps.
\item \textbf{Large restart period $M\to\infty$.} The schedule \eqref{eq:eta} approaches a simple cosine decay without reset; the analysis remains unchanged because only the average $\frac{1}{T}\sum\eta_t$ appears.
\item \textbf{Very small $M$ (restart each iteration).} Then $\eta_t\equiv\eta_{\max}$ and SGDR collapses to vanilla SGD with constant step size; the theorem recovers the classic bound.
\end{itemize}

\end{document}

% ===== CORRECTION SUMMARY =====
% 1. Step 6: Fixed the inequality by proving Lemma 2 and tightening the step‑size assumption to $\eta_{\max}\le 1/L$.
% 2. Step 7: Removed the erroneous expectation on the deterministic learning rate.
% 3. Step 9: Added Lemma 3 to rigorously bound the average learning rate for any horizon $T$, handling the case where $T$ is not a multiple of $M$.
% 4. Updated Assumption 5 to reflect the correct bound on $\eta_{\max}$ and introduced the auxiliary bound (A6).
% 5. Clarified the definition of the random index $\tau$ in Step 10.